\subsection{Quadratic canonical-gauge changes on a primitive mixed line}

\begin{proposition}[DP48: a quadratic gauge cannot improve the path slope]\label{DP48}
Use all hypotheses and notation of DP42 over the unramified coefficient field $K$, including the nonzero integral-function restriction $V(t)$. For $\lambda\in\mathcal O^\times$ and $c\in K$ put
\[
 \Theta_c(t)=V(t)\exp\bigl(\lambda\beta(ut,\mathbf vt)+ct^2\bigr)
        =\sum_{n\ge0}a_nt^n.
\]
If $c=0$ set $r_c=-\infty$; otherwise set
$r_c=(s-\vp(c))/2$. Write $R_c=\max\{w_k,r_c\}$.
Then
\[
 \nu_w(\Theta_c)=
 \begin{cases}
 -\infty,&0\le w<R_c,\\
 \nu_w(V)\in[0,\infty),&w\ge R_c,
 \end{cases}
\]
for rational $w$, and
\[
 \vp(n!a_n)\ge(s-R_c)n-\frac{s_p(n)}{p-1},\qquad
 \liminf_{n\to\infty}\frac{\vp(n!a_n)}n=s-R_c.
\]
If $c$ is integral, the sharp slope of DP42 is unchanged. If
$\vp(c)=-M\le-1$, the new slope is $(s-M)/2<0$, strictly worse.
\end{proposition}

\begin{proof}
The case $c=0$ is DP42. Otherwise $m=\vp(c)\in\Z$, by unramifiedness.
If $m\ge0$, then $r_c\le s/2<w_k$.
At every $w\ge w_k$, DP42 makes the primitive exponential integral, and $m+2w\ge s$ makes both $\exp(c\gamma^2T^2)$ and its inverse integral for $\vp(\gamma)=w$. Both have coefficient infimum zero. The prefactor is integral and nonzero, so Gauss multiplicativity gives the claimed finite value.
For $w<w_k$ sufficiently close to $w_k$, the quadratic exponential and its inverse remain integral, while DP42 makes $V\exp(\lambda\beta)$ unbounded. Their product cannot become bounded, since multiplication by the integral inverse would then contradict DP42. Monotonicity gives unboundedness at every smaller nonnegative weight. Thus $R_c=w_k$ in this case.

If $m=-M\le-1$, then, since $s\le1/2$,
$r_c=(s+M)/2>s>w_k$.
At every $w\ge r_c$ both exponential factors are integral with coefficient infimum zero, proving the finite equality again.
For $w<r_c$ sufficiently close to $r_c$, we have $w>w_k$ and
$0\le m+2w<s$.
The factor $V(\gamma T)\exp(\lambda\beta(u\gamma T,\mathbf v\gamma T))$ is nonzero and bounded by DP42.
The exponent $c\gamma^2T^2$ has initial form a nonzero scalar multiple of $T^2$, with nonzero derivative since $p$ is odd.
DP41's fraction-field multiplier lemma therefore proves unboundedness of the full product. Monotonicity extends it to every $w<r_c$.

At $w=R_c$ integrality gives $\vp(a_n)\ge-R_cn$.
If the coefficient liminf were greater than $-R_c$, a rational smaller weight sufficiently close to $R_c$ would have finite coefficient infimum, contradicting the established threshold. Hence $\liminf\vp(a_n)/n=-R_c$. Legendre's exact formula gives the two factorial conclusions.
\end{proof}

The source's difference of canonical logarithmic gauges is balanced, symmetric and biadditive; after logarithms it is bilinear, and after a linear logarithmic path and the $w_0$ diagonal it becomes $ct^2$. Thus DP48 applies once the prescribed expression is identified with the primitive mixed kernel up to an integral prefactor and this ambiguity. It does not supply that geometric identification. Nor does it replace the exact annihilation condition in the source's selected harmonic gauge-independence theorem. The integer valuation of $c$ separates the two cases; no statement for arbitrary ramified quadratic coefficients is asserted here.

